\subsection{Cell-এর আকার, Frequency Reuse ও Tower নির্বাচন}
পাহাড়, ভবন, অ্যান্টেনা ও শক্তির প্রভাবে বাস্তব সেলের সেবার পরিসর অনিয়মিত হয়। পরিকল্পনার চিত্রে ষড়ভুজ ব্যবহার করা হয়। ষড়ভুজ প্রায় বৃত্তাকার পরিসর বোঝায় এবং পাশাপাশি বসিয়ে ফাঁক ছাড়াই পুরো এলাকা দেখানো যায়। একটি বিশাল টাওয়ারের বদলে ছোট সেল ব্যবহার করলে ফ্রিকোয়েন্সি পুনর্ব্যবহার এবং বেশি ব্যবহারকারীকে সেবা দেওয়া সহজ হয়।

\textbf{ফ্রিকোয়েন্সি পুনর্ব্যবহার:} যথেষ্ট দূরের সেলে একই ফ্রিকোয়েন্সি আবার ব্যবহার করা যায়। কাছের সেলগুলোর বিঘ্ন নিয়ন্ত্রণ করতে পরিকল্পনা প্রয়োজন। মোবাইল সাধারণত সবচেয়ে উপযুক্ত সেবাদানকারী সেল বেছে নেয়। সেটি সবসময় সবচেয়ে কাছের টাওয়ার নয়। সংকেতের মান, কাজের চাপ ও নেটওয়ার্কের সিদ্ধান্তও নির্বাচনে ভূমিকা রাখে।

\textbf{GSM/CDMA-এর সেবার পরিসরের উদাহরণ:} প্রচলিত তুলনায় GSM-এর প্রায় ৩৫ km এবং CDMA-এর প্রায় ১১০ km সেলের পরিসর ধরা হয়। GSM-এর প্রচলিত Timing Advance-এর সীমা প্রথম সংখ্যাটির একটি প্রযুক্তিগত প্রেক্ষাপট। বিশেষ সম্প্রসারণে কার্যকর দূরত্ব ভিন্ন হতে পারে। CDMA-র সংখ্যাটি নির্দিষ্ট ঐতিহাসিক নকশার উদাহরণ; সব CDMA সেলের আবশ্যিক পরিসর নয়। ছোট শহুরে সেল বা ঘরের ভেতরের সেলও হতে পারে। শুধু প্রযুক্তির নাম দিয়ে বাস্তব সেলের ব্যাসার্ধ জানা যায় না।

\textbf{GSM-900-এর ফ্রিকোয়েন্সির হিসাব:} প্রচলিত GSM-900-এ ফোন থেকে বেস স্টেশনে Uplink হলো ৮৯০–৯১৫ MHz। বিপরীত পথে Downlink হলো ৯৩৫–৯৬০ MHz। প্রতিটি দিকের ব্যান্ডউইথ $915-890=25$ MHz। সমতুল্য Uplink ও Downlink ক্যারিয়ারের ব্যবধান ৪৫ MHz। রেডিও চ্যানেলের ব্যবধান ২০০ kHz এবং একটি TDMA ফ্রেমে ৮টি সময়ের ভাগ থাকে। GSM-এর অন্য ব্যান্ড ও সম্প্রসারিত ব্যান্ডে ফ্রিকোয়েন্সি আলাদা। এই সংখ্যা সব মোবাইল প্রজন্মের জন্য প্রযোজ্য নয়।

\subsection{একই মাধ্যম বহু গ্রাহক ভাগ করে কীভাবে?}
\textbf{GSM-এ Frequency Hopping:} নির্ধারিত ক্রম অনুসরণ করে একাধিক অনুমোদিত ক্যারিয়ার ফ্রিকোয়েন্সিতে পালাক্রমে সংকেত পাঠানোকে Frequency Hopping বলে। প্রেরক ও প্রাপক একই ক্রম অনুসরণ করে। শুধু বেশি ফ্রিকোয়েন্সি বেছে নেওয়াই এর কাজ নয়। নির্দিষ্ট ফ্রিকোয়েন্সির বিঘ্ন বা সংকেত ম্লান হওয়ার প্রভাব বিভিন্ন ফ্রিকোয়েন্সিতে ছড়িয়ে দিতে এই পদ্ধতি সাহায্য করতে পারে। Frequency Hopping এবং এক সেল থেকে অন্য সেলে Handoff পৃথক ঘটনা।

FDMA = Frequency Division Multiple Access; TDMA = Time Division Multiple Access; CDMA = Code Division Multiple Access। তিন পদ্ধতিতে ভাগ করার ভিত্তি যথাক্রমে ফ্রিকোয়েন্সি, সময় ও কোড।
\begin{tabularx}{\linewidth}{|p{.18\linewidth}|X|X|}
\hline
\textbf{পদ্ধতি} & \textbf{ভাগ করার ভিত্তি} & \textbf{সহজ উদাহরণ ও প্রজন্ম}\\\hline
FDMA & পৃথক ফ্রিকোয়েন্সি চ্যানেল & একই সময়ে পৃথক ফ্রিকোয়েন্সিতে যোগাযোগ; 1G-এর পরিচিত পদ্ধতি।\\\hline
TDMA & একই চ্যানেলে পৃথক সময়ের ভাগ & পালাক্রমে ছোট সময়ের অংশ বরাদ্দ; 2G GSM-এ ব্যবহৃত।\\\hline
CDMA & পৃথক কোড & কোড অনুযায়ী তথ্য আলাদা করা; 2G-এর কিছু প্রযুক্তি পরিবার ও 3G WCDMA।\\\hline
\end{tabularx}
একই মাধ্যম বহু ব্যবহারকারী ভাগ করে ব্যবহারের ধারণা হলো Multiple Access। TDMA-তে সময় ভাগ করা হলেও ব্যবহারকারীর পুরো টেলিফোন কথোপকথন হাফ-ডুপ্লেক্স হয়ে যায় না। সংযোগের দিক ও বেতার সম্পদ ভাগ করার পদ্ধতি আলাদা বিষয়।

\textbf{GSM ও CDMA-র পার্থক্য:} প্রচলিত GSM-এ প্রতিটি ফ্রিকোয়েন্সি চ্যানেলের মধ্যে পৃথক সময়ের ভাগ বরাদ্দ করা হয়। অর্থাৎ FDMA ও TDMA দুটিই কাজে লাগে। CDMA-তে পৃথক কোড দিয়ে একই যোগাযোগসম্পদে ব্যবহারকারীদের তথ্য আলাদা করা হয়। কোড ব্যবহারকারীর ব্যক্তিগত ফোন নম্বর নয়। এই ব্যবস্থায় উন্নত শক্তি নিয়ন্ত্রণ ও উপযুক্ত গ্রাহক যন্ত্র প্রয়োজন। সব CDMA ব্যবস্থার ব্যাটারির ব্যবহার, কার্যকর দূরত্ব বা কথার মান সমান নয়।
\chTwoPair{\textbf{Subscriber-এর পরিচয়:} SIM = Subscriber Identity Module। গ্রাহকের পরিচয় ও তা যাচাইয়ের তথ্য অপারেটরের সেবা নিতে সাহায্য করে। IMSI (International Mobile Subscriber Identity) গ্রাহকের স্থায়ী এবং TMSI (Temporary Mobile Subscriber Identity) অস্থায়ী পরিচয়। IMEI (International Mobile Equipment Identity) সাধারণত ১৫ অঙ্কের হ্যান্ডসেট পরিচয়। এটি গ্রাহকের পরিচয় বা ফোন নম্বর নয়।}{\textbf{Handset বদলানোর শর্ত:} প্রচলিত GSM-এ অন্য সমর্থিত হ্যান্ডসেটে SIM ব্যবহার করা যায়। বেতার ব্যান্ড, প্রযুক্তি, অপারেটরের লক ও অনুমতি সামঞ্জস্যপূর্ণ হতে হবে। কিছু CDMA হ্যান্ডসেট অপারেটর সক্রিয় করে দেয়। কিছু হ্যান্ডসেটে RUIM (Removable User Identity Module) থাকে। তাই সব CDMA ব্যবস্থায় কার্ড বা রোমিং নেই বলা ঠিক নয়।\par
\textbf{পুরোনো পরিচয়গুলোর নাম:} ESN (Electronic Serial Number) যন্ত্রের, MIN (Mobile Identification Number) গ্রাহকের এবং SID (System Identification) সিস্টেমের পরিচয় বোঝায়। ঐতিহাসিক সেলুলার বা CDMA ব্যবস্থা বুঝতে নামগুলো প্রয়োজন। সব নতুন নেটওয়ার্কে একই পরিচয়সূচক ব্যবহৃত হয় না।}

\subsection{প্রযুক্তির পূর্ণরূপ, চালুর সাল ও Rate}
\begin{multicols}{2}
\textbf{GSM, GPRS ও EDGE:} GSM: Global System for Mobile Communications। GPRS: General Packet Radio Service। EDGE: Enhanced Data rates for GSM Evolution। GPRS ও EDGE দিয়ে GSM-ভিত্তিক নেটওয়ার্কে উন্নত প্যাকেট ডেটা সেবা যুক্ত হয়।\par
\textbf{UMTS, LTE ও NR:} UMTS: Universal Mobile Telecommunications System। LTE: Long Term Evolution। NR: New Radio। পরীক্ষার উদ্দীপকে LTE দিয়ে সাধারণত 4G বোঝানো হয়। LTE-Advanced কঠোর IMT-Advanced লক্ষ্য পূরণের দিকে উন্নত রূপ।\par
\end{multicols}

\textbf{কয়েকটি নির্দিষ্ট চালুর সাল:} প্রথম 1G নেটওয়ার্কের ১৯৭৯/১৯৮৩ সালের প্রেক্ষাপট আগের 1G অংশে আছে। Finland-এ Nokia-র যন্ত্রে প্রথম আনুষ্ঠানিক GSM কল হয় ১ জুলাই ১৯৯১ সালে। TeliaSonera-র Stockholm LTE সেবায় ২০০৯ সালের ডিসেম্বরে Ericsson-এর নেটওয়ার্ক ব্যবহৃত হয়। সেবাটি বাজারে 4G নামে পরিচিত। একটি দেশে সেবা উদ্বোধনের সাল পুরো বিশ্বের প্রজন্ম বদলের সাল নয়।

\textbf{5G-এর চালুর সাল:} ২০১৯ সালেই প্রাথমিক বাণিজ্যিক 5G সেবা চালু ছিল। IMT-2020 নামের মধ্যে ২০২০ থাকলেও সব দেশে প্রথম 5G চালুর সাল ২০২০ নয়। প্রযুক্তির উন্নয়ন, মান অনুমোদন ও স্থানীয় সেবা চালুর তারিখ পৃথক হতে পারে।

\textbf{FOMA:} Freedom Of Mobile multimedia Access হলো NTT DOCOMO-র ২০০১ সালে চালু করা 3G সেবার ঐতিহাসিক নাম। এটি পুরো 3G প্রজন্মের পূর্ণরূপ নয়। WCDMA-তে W দিয়ে Wideband বোঝায়।

\textbf{মান ও প্রযুক্তির নাম:} IMT = International Mobile Telecommunications। ITU-র IMT-2000/Advanced/2020 যথাক্রমে 3G/4G/5G পরিচয়ের সঙ্গে সম্পর্কিত মানের পরিবার। 3GPP = Third Generation Partnership Project। এটি মান উন্নয়নের অংশীদারত্ব; ফোন বা একক বেতার লিংকের নাম নয়। HSPA = High Speed Packet Access। HSDPA (High Speed Downlink Packet Access) ডাউনলোড এবং HSUPA (High Speed Uplink Packet Access) আপলোডের উন্নয়নের পরিচিত নাম। RAT = Radio Access Technology; GSM বা LTE-এর মতো বেতার সংযোগের প্রযুক্তি বোঝায়।

\textbf{3G-এর প্রযুক্তি ও 4G-এর সাবক্যারিয়ার:} 3G-তে UMTS/WCDMA ছাড়াও CDMA2000 পরিবার আছে। TD-CDMA ও TD-SCDMA (Time Division Synchronous CDMA)-ও এই প্রজন্মের প্রযুক্তি। একটি নাম পুরো 3G-এর একমাত্র পদ্ধতি নয়। OFDM (Orthogonal Frequency Division Multiplexing)-এ ডেটা সমন্বিত বহু সাবক্যারিয়ারে ভাগ করে বহন করা হয়। OFDMA-তে সাবক্যারিয়ারের সম্পদ বিভিন্ন ব্যবহারকারীকে বরাদ্দ দেওয়া হয়। LTE-এর Downlink এবং Wi-Fi/WiMAX-এর উপযুক্ত সংস্করণে এসব ব্যবহৃত হয়। এটি পৃথক তামার তারে Parallel Transmission-এর একই ধারণা নয়।

\textbf{3G-এর প্রাথমিক হারের উদাহরণ:} পুরোনো IMT-2000 পাঠে দ্রুত চলমান অবস্থায় প্রায় ১৪৪ kbps, হাঁটা বা ধীর চলাচলে ৩৮৪ kbps এবং স্থির বা ঘরের ভেতরে ২ Mbps-এর লক্ষ্য পাওয়া যায়। পরে HSPA-তে ডেটা হার বাড়ে। তাই 3G-এর গতি সবসময় শুধু ২ Mbps নয়। এই মানগুলো ব্যবহারে চলাচল, বেতার প্রযুক্তি ও পরিবেশের শর্ত বিবেচনা করতে হয়।

\textbf{ডেটা হারের শর্ত:} 4G-তে ১০০ Mbps (দ্রুত চলাচল) ও ১ Gbps (ধীর চলাচল) পরিচিত গবেষণা লক্ষ্য। IMT-2020/5G-এর আদর্শ সর্বোচ্চ Downlink ২০ Gbps এবং Uplink ১০ Gbps। প্রতি বর্গকিলোমিটারে প্রায় ১০ লক্ষ যন্ত্র সংযোগের লক্ষ্যও আছে। এগুলো নির্দিষ্ট পরীক্ষার শর্তের লক্ষ্য। প্রত্যেক ফোনের নিশ্চিত গতি বা সব লক্ষ্য একই সঙ্গে পাওয়ার প্রতিশ্রুতি নয়। পরিষেবার প্রচলিত হারের ছক অংশ \ref{ch2:c01-12}-এ আছে।

\textbf{WWWW-এর পরিচয়:} World Wide Wireless Web দিয়ে 5G-এর বিস্তৃত বেতার সেবার ধারণা বোঝানো হয়। Wireless World Wide Web শব্দক্রমও প্রচলিত। এটি IMT-2020 বা NR-এর মতো আনুষ্ঠানিক বেতার মান নয়। IoT, কম বিলম্ব, বেশি বহনক্ষমতা ও 4K ভিডিওর বৈশিষ্ট্যগুলো একত্রে 5G-এর প্রসঙ্গ নির্দেশ করতে পারে। শুধু 4K ভিডিও অন্য উপযুক্ত ব্রডব্যান্ডেও চলে। তাই ভিডিওর রেজল্যুশন একা প্রজন্ম শনাক্ত করে না।

\textbf{আর্কিটেকচারের সীমা:} BTS–BSC–MSC-সহ চিত্রটি প্রচলিত GSM/2G-এর সরল মডেল। 4G/5G-এর বেতার অংশ ও মূল নেটওয়ার্কে নাম এবং দায়িত্ব ভিন্ন হয়। নতুন প্রযুক্তির সম্ভাবনা ব্যাখ্যা করা যায়। তবে শুধু 5G থাকলেই দূর থেকে অস্ত্রোপচার বা স্বয়ংচালিত গাড়ির সব ব্যবস্থা নিরাপদে চলবে—এমন সিদ্ধান্ত দেওয়া যায় না।

\begin{keypointbox}{MIMO: একাধিক Antenna কেন?}
MIMO (Multiple Input Multiple Output)-তে প্রেরক ও প্রাপক প্রান্তে একাধিক অ্যান্টেনা ব্যবহার করা হয়। উপযুক্ত পরিবেশে একই সময় ও ফ্রিকোয়েন্সির সম্পদে একাধিক ডেটা প্রবাহ বহন করা যায়। এতে সংযোগের বহনক্ষমতা বাড়তে পারে। LTE/4G ও আধুনিক Wi-Fi-তে এই ধারণা ব্যবহৃত হয়। রাস্তার একাধিক লেন দিয়ে পাশাপাশি গাড়ি যাওয়ার সঙ্গে এর তুলনা করা যায়। তবে শুধু অ্যান্টেনা বাড়ালেই সব পরিবেশে একই গতি পাওয়া যায় না। সংকেত, চ্যানেল ও যন্ত্রের সক্ষমতাও প্রয়োজন।
\end{keypointbox}

